% FILE: 02_lineage_and_context.tex
% Project: adversarial
% Thesis Role: adversarial-testing-and-hostile-assumption-verification

\section{Lineage and Context}
\label{sec:adversarial-lineage}

\subsection{2016 Language-Model Lesson}
\label{sec:adversarial-lineage-2016}

The thesis lineage begins with a 2016 language-model experiment whose conclusion was that local
text imitation does not conserve consequence. A model that can reproduce the surface form of a
process description does not thereby understand the process, cannot detect when the process
deviates from its declared form, and cannot produce evidence that a process claim is true. This
lesson is the deepest motivation for the adversarial corpus.

The adversarial corpus operationalizes this lesson as a falsification requirement: it is
insufficient to claim that a process intelligence system produces conforming output. The claim
must survive an adversary who controls the input. A system that accepts timestamp-laundered
logs, event-deleted traces, or forged signatures has not achieved process intelligence; it has
achieved a more sophisticated form of text imitation. The ten scenarios in
\texttt{simulate\_laundering\_spoofing.py} are the direct descendants of the 2016 lesson,
translated into executable Python tests that assert the system cannot be fooled.

\subsection{Enterprise Process Gap}
\label{sec:adversarial-lineage-enterprise}

The enterprise process gap---the observation that declared processes and actual runtime
processes routinely diverge---is the empirical ground for the adversarial corpus. The gap
manifests in three forms identified by the adversarial canon
(\texttt{sources/papers/adversarial-canon.md}): Sybil attacks on event correlation (forged
case identifiers causing state-space explosion), non-deterministic reality receipts
(out-of-band state changes misinterpreted as noise by discovery algorithms), and the XES
flat-earth problem (lossy ETL transformations that allow adversarial event selection without
schema violation).

These forms of the enterprise process gap are not theoretical. They are the adversarial
surfaces that any process intelligence system deployed in an ERP, CRM, or SCM environment
will face. The adversarial corpus provides the rejection mechanisms: hash chain integrity
verification detects event deletion, signature authenticity checks detect timestamp
modification and key forgery, and the \texttt{reject\_unhashed\_dataframe} rule blocks the
entire class of raw mutable-data injection attacks.

\subsection{Relationship to the Chatman Equation}
\label{sec:adversarial-lineage-chatman-equation}

The Chatman Equation $A = \mu(O^*)$ asserts that accountable process output is a function of
a specified object-centric transformation. The adversarial corpus provides the negative
complement of this equation: it defines what happens when $\mu$ is forged, when the inputs to
$O^*$ are corrupted, and when the receipt certifying $A$ is manufactured from an adversarial
key.

The Chatman Constraint (\texttt{sources/wasm4pm-compat/adversarial-type-law.md}, Law I) is
the injectivity requirement that makes the equation non-trivially falsifiable:
$\forall S_a, S_b \in \text{StateSpace},\ S_a \neq S_b \implies C(S_a) \neq C(S_b)$.
If this constraint is not enforced, state-space aliasing allows an adversary to forge a
computationally indistinguishable but semantically distinct witness state. The adversarial
corpus demonstrates that the Chatman Constraint must be a structural law at the
compiler/scaffolding level, not a runtime check, because a runtime check can itself be forged
if the witness lattice is malformed.

\subsection{Relationship to Receipts and Replay}
\label{sec:adversarial-lineage-receipts}

The receipt mechanism in the process-intelligence foundry is the primary defense against
adversarial process claims. The ADVERSARIAL\_RECEIPT in \texttt{receipts/RECEIPT\_REGISTRY.md}
certifies the adversarial workflow under the Chicago TDD hostile assumptions doctrine,
witnessing that challenges cover raw laundering, metric fabrication, and conformance theater.

Replay is the adversarial surface that the receipt mechanism must foreclose. The adversarial
audit report (\texttt{audits/adversarial\_audit\_v30.1.1.md}) identifies memory-snapshot
replay as Vulnerability Vector 2: snapshotting WASM linear memory and replaying evidence
blocks in different execution contexts. The mitigation is monotonic session-based epoch
counters and hash-salted state transitions, ensuring that replayed sessions and signature
reuse across epochs are rejected at the boundary. Witness Lattice Monotonicity
($W_{\text{new}} = W_{\text{old}} \sqcup w_{\text{step}}$) provides the same protection at
the type-law level: any non-monotonic transition triggers a self-halt before a fraudulent
receipt can be emitted.
